\section{Por qué las trayectorias de razonamiento parecen humanas: emergencia, búsqueda y la frontera de lo verificable}

Con tokens intermedios y outcome-guided training, el modelo puede producir procesos de resolución largos y estructurados. En clase se le llamó la «beauty» del LLM reasoning: sin programar explícitamente un buscador exhaustivo, la token-to-token generation también puede mostrar descomposición, estimación, retroceso y verificación local. Este fenómeno merece estudiarse, pero la interpretación errónea más peligrosa es tomar la superficie lingüística como si fuera identidad cognitiva.

\subsection{La estructura que emerge de la generación de tokens}

La slide 27 contrasta los LLM modernos con la exhaustive search de la inteligencia artificial clásica. Una interpretación más precisa es esta: el modelo comprime una gran cantidad de patrones en sus parámetros y, al generar, avanza a lo largo de una trayectoria o de unas pocas; no necesariamente enumera de forma explícita el árbol de estados completo. La ventaja es que la ruta de razonamiento es compacta y puede combinarse con el lenguaje natural y con interfaces de herramientas; el costo es una cobertura incompleta, y las ramas erróneas no necesariamente se detectan de forma sistemática.

\lecturefigure{slide-27.jpg}{El atractivo del LLM reasoning: un proceso estructurado puede surgir de la generación de tokens}{27}

\teachervoice{00:31:20--00:32:10, el docente usa como contraste la valoración de Kasparov sobre Deep Blue y enfatiza que las trayectorias que vemos ahora ya no son solo programas de búsqueda escritos a mano. Esta nota advierte: se trata de una descripción de la forma del comportamiento; no basta para demostrar que el modelo no realiza una búsqueda implícita, ni demuestra que la trayectoria sea fiel.}

\begin{importantbox}{Cómo debe entenderse aquí la «emergencia»}
La emergencia no es una etiqueta mágica; significa que el objetivo de entrenamiento no codificó paso por paso el guion «primero estimar la escala, luego buscar una descomposición y al final verificar», y aun así el modelo, con la escala, los datos y el entrenamiento posterior, presenta esa estructura reutilizable. La investigación debe seguir preguntando qué datos y recompensas la inducen, en qué distribuciones falla y si un verifier puede revisarla.
\end{importantbox}

\subsection{El problema aritmético de 2025: qué aporta una trayectoria larga}

El problema pide usar cada número del 1 al 10 una sola vez, solo con sumas y multiplicaciones, para obtener 2025. Una construcción concisa es
\[
(10\times4+5)\times(9\times3+8+7+2+1)=45\times45=2025.
\]
El primer grupo forma 45 y el segundo también forma 45. La slide 28 contrasta la solución concisa de estilo humano con la trayectoria larga del thinking mode de Gemini 2.0 en diciembre de 2024: el modelo primero nota que 2025 es $45^2$, luego transforma el objetivo en construir dos 45 y, tras una exploración muy larga de candidatos, da la respuesta.

\lecturefigure{slide-28.jpg}{El problema aritmético de 2025: una trayectoria de razonamiento larga que va de la estimación de magnitud a la construcción de dos 45}{28}

Al leer la figura, primero se observa la \textbf{reformulación del problema} en la trayectoria de la derecha: un objetivo grande sugiere multiplicación y la estructura de cuadrado sugiere dos factores; después se revisa si el proceso de búsqueda mantiene la restricción «cada número se usa una sola vez»; solo al final se mira la respuesta. El valor de un CoT largo no está en el número de palabras, sino en hacer explícitas las restricciones, las heurísticas y los objetivos intermedios, de modo que un verifier pueda revisarlos paso a paso y que los tokens posteriores lean la estructura de $45$ ya descubierta.

\begin{warningbox}{Tres tipos de alucinación en las trayectorias largas}
Primero, adivinar la respuesta y después fabricar un proceso plausible; segundo, violar restricciones a mitad del camino mientras la expresión final resulta correcta por casualidad; tercero, exploración repetitiva, circular o inútil que consume una gran cantidad de tokens. La evaluación debe revisar a la vez la expresión final, el conjunto de números usados, las restricciones de operación y el costo de la trayectoria.
\end{warningbox}

\subsection{Aprender el proceso de descubrimiento, no solo cargar los resultados descubiertos}

La cita de Rich Sutton enfatiza que se busca que el agent aprenda el discovery process, no que solo contenga las respuestas que los humanos ya descubrieron. Para un reasoning model, esto significa que los datos de entrenamiento no pueden consistir solo en versiones finales de alta calidad; también deben ofrecer una distribución de problemas que pueda explorarse, retroalimentación confiable y oportunidades para recuperarse de los fallos. De lo contrario, el modelo se parece más a un enorme recuperador de soluciones y, ante combinaciones nuevas, sigue careciendo de un mecanismo de descubrimiento.

\lecturefigure{slide-29.jpg}{La perspectiva de la Bitter Lesson: el objetivo es aprender el proceso de descubrimiento}{29}

\teachervoice{00:32:10--00:39:00, el docente usa la trayectoria aritmética larga para mostrar que el modelo forma heurísticas parecidas a las humanas, pero recuerda repetidamente que un LLM es un probabilistic model. Nota de clase: los pasos legibles son un activo de ingeniería que sirve para la verificación y el entrenamiento; «parecer humano» es solo una descripción y no debe tomarse como evidencia de corrección.}

Aquí la búsqueda no desaparece, sino que se divide en tres niveles: la compresión de patrones en los parámetros, la exploración local en la secuencia de tokens y, de forma opcional, la symbolic search/tool use externa. En la Q\&A, el docente aclaró además que el modelo puede escribir Python o invocar un buscador; eso corresponde al uso de herramientas. Al estudiar la «reasoning ability básica» puede examinarse por separado la generación sin búsqueda externa, pero los sistemas de producto normalmente deben combinar ambas cosas, en lugar de plantear search y learning como una oposición ideológica.

\subsection{Los límites de capacidad del RL reasoning: no todas las tareas tienen un verifier automático}

La slide 30 presenta las ventajas y desventajas centrales del RL finetuning: generaliza bien en tareas cuyas respuestas pueden verificarse automáticamente, pero muchas tareas abiertas carecen de una función de evaluación única, barata y confiable. El exact match en matemáticas, las pruebas de código y el resultado de una partida de ajedrez pertenecen a zonas relativamente favorables; las recomendaciones de política pública, la investigación de textos largos, el diseño de productos y el juicio social suelen implicar múltiples objetivos, retroalimentación diferida y conflictos de valores.

\lecturefigure{slide-30.jpg}{Los límites del RL finetuning: tareas verificables automáticamente frente a tareas abiertas}{30}

Las tareas pueden clasificarse en cuatro tipos según las condiciones del verifier:
\begin{center}
\begin{tabularx}{0.96\textwidth}{p{0.19\textwidth}X X}
\toprule
Tipo & Retroalimentación típica & Riesgo principal \\
\midrule
Respuesta única & exact match, tolerancia numérica & errores de análisis sintáctico, filtración de la respuesta \\
Resultado ejecutable & pruebas unitarias, recompensa del simulador & pruebas incompletas, reward hacking \\
Preferencia por pares & elección de un judge humano o de un modelo & sesgo del judge, preferencia de estilo \\
Objetivo abierto a largo plazo & resultados de negocio, impacto social & atribución diferida, conflicto entre objetivos, costos irreversibles \\
\bottomrule
\end{tabularx}
\end{center}

\begin{knowledgebox}{Verifier frontier}
La siguiente etapa del entrenamiento de razonamiento no consiste en trasladar mecánicamente las recetas de RL existentes a todos los dominios, sino en construir retroalimentación de proceso más confiable, entornos ejecutables, pruebas contrafactuales y criterios de aceptación con múltiples métricas. Sin esta infraestructura, la «recompensa» en tareas abiertas puede fácilmente medir solo el formato, la complacencia o indicadores sustitutos de corto plazo.
\end{knowledgebox}

\subsection{Por qué el siguiente paso es la agregación y la recuperación}

Aunque una sola reasoning path ya sea muy capaz, la generación sigue siendo un proceso aleatorio: una misma pregunta puede producir distintos pasos intermedios y respuestas. La slide 31 anticipa dos direcciones complementarias: la aggregation usa varios candidatos para estimar la probabilidad de la respuesta, y la retrieval complementa con conocimiento pertinente procedente del exterior o del contexto. La slide 32 recuerda de nuevo que un LLM es un next-token probability model, no un ser humano; el flujo natural de «generar primero la trayectoria y luego dar la respuesta» puede seguir desalineado respecto del objetivo probabilístico.

\lecturefigure{slide-31.jpg}{Dos vías para mejorar aún más: Aggregation y Retrieval}{31}

\lecturefigure{slide-32.jpg}{Incluso con reasoning, la next-token decoding puede no coincidir con el objetivo de la respuesta final}{32}

\teachervoice{00:39:00--00:40:17, el docente pasa de las limitaciones del verifier en RL a las mejoras en inference-time: cuando una sola trayectoria es inestable, no hay que suponer que basta con una generación; conviene aprovechar la estructura estadística de varias muestras o recuperar primero el conocimiento pertinente y después resolver el problema.}

\subsection{Resumen de la sección}

Una reasoning trace larga puede contener objetivos intermedios, restricciones y verificación local, pero su corrección no puede demostrarse por el hecho de que «parezca humana». El RL tiene ventaja en tareas verificables automáticamente; el cuello de botella para extenderlo es el verifier de los objetivos abiertos. Lo que sigue es resolver la discrepancia entre una sola trayectoria aleatoria y la probabilidad de la respuesta final.
